\expandafter\ifx\csname CAPETINCLUDEMODE\endcsname\relax
\documentclass[11pt,a4paper]{article}
% preamble.tex — packages et commandes communes a toutes les fiches CAPET SII IE
\input{glyphtounicode}
\pdfgentounicode=1
\usepackage{cmap}
\usepackage[T1]{fontenc}
\usepackage[utf8]{inputenc}
\usepackage[french]{babel}
\usepackage[a4paper, top=1.6cm, bottom=1.6cm, left=1.4cm, right=1.4cm]{geometry}
\setlength{\headheight}{14.5pt}
\usepackage{amsmath, amssymb, amsfonts}
\usepackage{siunitx}
\sisetup{locale=FR, output-decimal-marker={,}, inter-unit-product=\ensuremath{{}\cdot{}}}
\usepackage{array, booktabs, multirow, tabularx}
\usepackage{xcolor}
\usepackage[most, breakable]{tcolorbox}
\usepackage{multicol}
\usepackage{enumitem}
\usepackage{mathtools}
\usepackage{lmodern}
\usepackage{microtype}
\usepackage{fancyhdr}
\usepackage{tikz}
\usetikzlibrary{positioning,arrows.meta}
\usepackage{lastpage}

% ── Couleurs par domaine ──────────────────────────────────────────────
\definecolor{colElec}{RGB}{26, 82, 148}       % Électrotechnique — bleu foncé
\definecolor{colElecLight}{RGB}{210, 225, 245} % fond bleu clair
\definecolor{colAuto}{RGB}{34, 120, 66}        % Automatique — vert
\definecolor{colMeca}{RGB}{185, 90, 20}        % Mécanique — orange
\definecolor{colTherm}{RGB}{160, 40, 40}       % Thermique — rouge
\definecolor{colQT}{RGB}{232, 244, 233}        % fond question type (vert très léger)
\definecolor{colJury}{RGB}{255, 243, 205}      % fond jury (jaune ambre)
\definecolor{colJuryFrame}{RGB}{200, 130, 0}   % bordure jury
\definecolor{colGrayLight}{RGB}{245, 245, 247} % fond gris neutre
\definecolor{colSeparator}{RGB}{180, 180, 180} % gris séparateur

% ── Style tcolorbox de section principale ────────────────────────────
\tcbset{
  ficheSection/.style={
    breakable,
    enhanced,
    colback=#1!7!white,
    colframe=#1,
    colbacktitle=#1,
    coltitle=white,
    fonttitle=\bfseries\normalsize,
    left=4pt, right=4pt, top=3pt, bottom=4pt,
    before skip=6pt, after skip=6pt,
  }
}

% ── Style encadré question type ──────────────────────────────────────
\tcbset{
  questiontype/.style={
    breakable,
    enhanced,
    colback=colQT,
    colframe=colElec!55!black,
    colbacktitle=colElec!55!black,
    coltitle=white,
    fonttitle=\bfseries\small,
    left=4pt, right=4pt, top=3pt, bottom=3pt,
    before skip=5pt, after skip=5pt,
  }
}

% ── Style encadré jury ───────────────────────────────────────────────
\tcbset{
  jury/.style={
    breakable,
    enhanced,
    colback=colJury,
    colframe=colJuryFrame,
    fonttitle=\bfseries\normalsize,
    left=4pt, right=4pt, top=3pt, bottom=3pt,
  }
}

% ── Mise en page ─────────────────────────────────────────────────────
\pagestyle{fancy}
\fancyhf{}
\renewcommand{\headrulewidth}{0.4pt}
\renewcommand{\footrulewidth}{0.3pt}

\setlist[itemize]{leftmargin=1.2em, itemsep=1pt, topsep=2pt, parsep=0pt}
\setlist[enumerate]{leftmargin=1.2em, itemsep=3pt, topsep=2pt, parsep=0pt}

\setlength{\parskip}{2pt}
\setlength{\columnsep}{10pt}

\def\CAPETBEGINFILE{\begin{document}}
\def\CAPETENDFILE{\end{document}}
\else
\def\CAPETBEGINFILE{}
\def\CAPETENDFILE{}
\fi

\fancyhead[L]{\textbf{\textcolor{colTherm}{T05}} - Phenomenes thermiques}
\fancyhead[R]{\small\textit{CAPET SII IE - Session 2026}}
\fancyfoot[C]{\small\thepage\,/\,\pageref{LastPage}}

\CAPETBEGINFILE

% -----------------------------------------------------------------------
%  BLOC EN-TeTE
% -----------------------------------------------------------------------

\begin{tcolorbox}[
  enhanced, breakable,
  colback=colTherm,
  colframe=colTherm,
  left=6pt, right=6pt, top=5pt, bottom=5pt,
  before skip=0pt, after skip=4pt
]
  {\color{white}\bfseries\large T05 \quad Phenomenes thermiques}\\[2pt]
  {\color{colTherm!30!white}\small
    \textbf{Domaine :} Thermique / Instrumentation \quad
    \textbf{Presence :} \textbf{8/13} sujets (2013-2026) - theme \emph{frequent}
  }\\[1pt]
  {\color{colTherm!20!white}\footnotesize
    Sous-themes : couple thermique equivalent - resistance thermique - capteurs (PT100, thermocouple) - transmetteur 4-20\,mA - vapeur saturee
  }
\end{tcolorbox}

% -----------------------------------------------------------------------
%  1. RAPPELS DE COURS
% -----------------------------------------------------------------------

\begin{tcolorbox}[
  ficheSection=colTherm,
  title={1. Rappels de cours},
]

% -- 1.1 Couple thermique equivalent ----------------------------------
\textbf{1.1 Couple thermique equivalent sur un cycle}

\smallskip
Un moteur soumis - un cycle de couples $C_{m1}, C_{m2}, \ldots$ de durees $\Delta t_1, \Delta t_2, \ldots$ (periode totale $T$) est
caracterise par son couple thermique equivalent :
\[
  C_{th} = \sqrt{\frac{\sum_{i} C_{mi}^2 \cdot \Delta t_i}{T}}
\]
Le courant thermique equivalent $I_{th}$ s'obtient de meme ($I_{th}^2 = \sum I_i^2\Delta t_i/T$).
$C_{th}$ doit etre inferieur au couple nominal du moteur pour garantir l'echauffement admissible.

\medskip
\hrule\vspace{4pt}

% -- 1.2 Resistance thermique -----------------------------------------
\textbf{1.2 Resistance thermique et modele thermique}

\smallskip
\begin{center}
\renewcommand{\arraystretch}{1.45}
\begin{tabular}{llll}
\toprule
\textbf{Grandeur} & \textbf{Expression} & \textbf{Unite} & \textbf{Analogie electrique} \\
\midrule
Flux thermique (conduction) & $\Phi = \lambda S\,\Delta T/e$ & W & courant \\[4pt]
Resistance thermique (conduction) & $R_{th} = e/(\lambda S)$ & K/W & resistance \\[4pt]
Resistance thermique (convection) & $R_{th,conv} = 1/(hS)$ & K/W & resistance \\[4pt]
elevation de temperature & $\Delta T = R_{th} \cdot P$ & K & loi d'Ohm \\[4pt]
Temperature jonction & $T_j = T_{amb} + R_{th,ja} \cdot P_{dissipee}$ & $^\circ\mathrm{C}$ & \\[4pt]
Constante de temps & $\tau_{th} = R_{th} \cdot C_{th,cap}$ & s & $\tau = RC$ \\
\bottomrule
\end{tabular}
\end{center}

\medskip
\hrule\vspace{4pt}

% -- 1.3 Capteurs de temperature --------------------------------------
\textbf{1.3 Capteurs de temperature}

\smallskip
\begin{center}
\renewcommand{\arraystretch}{1.35}
\begin{tabular}{lllll}
\toprule
\textbf{Capteur} & \textbf{Principe} & \textbf{Plage} & \textbf{Sortie} & \textbf{Precision} \\
\midrule
PT100 & Resistif ($R$ croet avec $T$) & $-200$ - $+850^\circ\mathrm{C}$ & $\Omega$ & $\pm 0{,}3^\circ\mathrm{C}$ \\[3pt]
Thermocouple J & Seebeck (Fe/Constantan) & $-30$ - $+700^\circ\mathrm{C}$ & $\mu$V/K & $\pm1{,}5^\circ\mathrm{C}$ cl.1 \\[3pt]
Thermocouple K & Seebeck (NiCr/Ni) & $-200$ - $+1260^\circ\mathrm{C}$ & $\mu$V/K & $\pm 2{,}2^\circ\mathrm{C}$ cl.1 \\[3pt]
DS1820 & Numerique, bus 1-Wire & $-55$ - $+125^\circ\mathrm{C}$ & bits & $\pm0{,}5^\circ\mathrm{C}$ \\
\bottomrule
\end{tabular}
\end{center}

\smallskip
PT100 : $R(T) = 100\,(1 + \alpha T)$ avec $\alpha = 3{,}85\times10^{-3}\,\text{K}^{-1}$.
Thermocouple : $e_{AB} = k_S(T_1 - T_2)$ (effet Seebeck). \textbf{Toujours compenser la jonction froide} si $T_{ref} \neq 0^\circ\mathrm{C}$.

\medskip
\hrule\vspace{4pt}

% -- 1.4 Transmetteur 4-20 mA -----------------------------------------
\textbf{1.4 Transmetteur 4-20 mA}

\smallskip
Relation lineaire entre grandeur physique $T$ et courant $I_{out}$ :
\[
  I_{out} = 4 + 16 \cdot \frac{T - T_{min}}{T_{max} - T_{min}} \quad [\text{mA}]
\]
$4\,\text{mA} \leftrightarrow T_{min}$ ; $20\,\text{mA} \leftrightarrow T_{max}$ ; $\Delta I = 16\,\text{mA}$ sur toute la plage.

\end{tcolorbox}

% -----------------------------------------------------------------------
%  2. FORMULAIRE SYNTHeTIQUE
% -----------------------------------------------------------------------

\begin{tcolorbox}[
  ficheSection=colTherm,
  title={2. Formulaire synthetique},
]
\small
\begin{multicols}{2}
\renewcommand{\arraystretch}{1.45}
\begin{tabular}{@{}p{4.2cm}ll@{}}
\toprule
\textbf{Formule} & \textbf{Usage} & \textbf{U.} \\
\midrule
\multicolumn{3}{@{}l}{\footnotesize\textit{Thermique moteur}} \\
$C_{th}^2 = \sum C_{mi}^2\Delta t_i\,/\,T$ & Couple therm. & Nem \\
$I_{th}^2 = \sum I_i^2\Delta t_i\,/\,T$ & Courant therm. & A \\
\midrule
\multicolumn{3}{@{}l}{\footnotesize\textit{Resistance thermique}} \\
$R_{th} = e/(\lambda S)$ & Conduction & K/W \\
$R_{th} = 1/(hS)$ & Convection & K/W \\
$\Delta T = R_{th} \cdot P$ & echauffement & K \\
$T_j = T_{amb} + R_{th,ja} \cdot P$ & Jonction & $^\circ\mathrm{C}$ \\
\bottomrule
\end{tabular}

\columnbreak

\renewcommand{\arraystretch}{1.45}
\begin{tabular}{@{}p{4.2cm}ll@{}}
\toprule
\textbf{Formule} & \textbf{Usage} & \textbf{U.} \\
\midrule
\multicolumn{3}{@{}l}{\footnotesize\textit{Capteurs}} \\
$R(T) = 100(1+\alpha T)$ & PT100 & $\Omega$ \\
$e = k_S(T_1 - T_2)$ & Thermocouple & V \\
\midrule
\multicolumn{3}{@{}l}{\footnotesize\textit{Transmetteur 4-20\,mA}} \\
$I_{out} = 4 + 16\cdot(T-T_{min})/\Delta T_m$ & Conversion & mA \\
$\Delta T_m = T_{max} - T_{min}$ & Plage capteur & $^\circ\mathrm{C}$ \\
\midrule
\multicolumn{3}{@{}l}{\footnotesize\textit{Effet Joule}} \\
$P_J = R\,I^2 = U\,I$ & Source chaleur & W \\
$\tau_{th} = R_{th} \cdot C_{th,cap}$ & Cste de temps & s \\
\bottomrule
\end{tabular}
\end{multicols}
\normalsize
\end{tcolorbox}

% -----------------------------------------------------------------------
%  3. QUESTIONS TYPES
% -----------------------------------------------------------------------

\begin{tcolorbox}[
  ficheSection=colTherm,
  title={3. Questions types - extraites des sujets reels},
]

% -- QT1 -------------------------------------------------------------

% Mini-procedure de resolution type copie concours
\begin{tcolorbox}[questiontype,title={Procedure type copie concours (5 etapes)}]
\begin{enumerate}[leftmargin=1.5em]
  \item \textbf{Qualifier le modele} : hypotheses, domaine de validite, grandeurs d'entree/sortie.
  \item \textbf{Ecrire les lois} : equations de base avec symboles et unites SI.
  \item \textbf{Isoler l'inconnue} : expression litterale, puis application numerique tracee.
  \item \textbf{Controler l'ordre de grandeur} : coherence physique et unite finale.
  \item \textbf{Conclure concours} : phrase courte orientee decision technique.
\end{enumerate}
\end{tcolorbox}

% Visuel de methode : chemin de resolution CAPET
\begin{center}
\begin{tikzpicture}[font=\small, node distance=7mm and 6mm, >=Stealth]
  \node[draw, rounded corners, fill=colElec!8, inner sep=4pt] (n1) {Notions};
  \node[draw, rounded corners, fill=colElec!8, right=of n1, inner sep=4pt] (n2) {Methode};
  \node[draw, rounded corners, fill=colElec!8, right=of n2, inner sep=4pt] (n3) {Application annale};
  \node[draw, rounded corners, fill=colElec!8, right=of n3, inner sep=4pt] (n4) {Verification jury};
  \draw[->, thick, colElec] (n1) -- (n2);
  \draw[->, thick, colElec] (n2) -- (n3);
  \draw[->, thick, colElec] (n3) -- (n4);
\end{tikzpicture}
\end{center}
\begin{tcolorbox}[questiontype,
  title={QT1 - Couple thermique equivalent sur un cycle \hfill\normalfont\footnotesize E2 sujet 2019}
]
\textbf{enonce.}
Un moteur subit un cycle de periode $T = 2048\,\text{s}$.
Les phases 1 et 2 sont caracterisees par $F_{e1} = 14{,}9\,\text{kN}$, $\Delta t_1 = 42\,\text{s}$ et $F_{e2} = 9{,}2\,\text{kN}$, $\Delta t_2 = 51\,\text{s}$.
Pour les phases 3 - 18 : $\sum C_{mi}^2\Delta t_i = 558\times10^6\,\text{N}^2\text{m}^2\text{s}$.
Calculer le couple thermique equivalent $C_{th}$.

\tcbline

\textbf{Corrige.}

\textbf{etape 1 - contribution des phases 1 et 2.}

(Les couples $C_{m1}$, $C_{m2}$ sont calcules - partir des forces et du bras de levier fournis en DT.)
\[
  \sum_\text{phases 1 et 2} C_{mi}^2\Delta t_i = C_{m1}^2 \cdot 42 + C_{m2}^2 \cdot 51
\]

\textbf{etape 2 - couple thermique total.}
\[
  C_{th} = \sqrt{\frac{C_{m1}^2\cdot42 + C_{m2}^2\cdot51 + 558\times10^6}{2048}}
\]
\textit{Erreur frequente : oublier de prendre la racine carree - la fin (on calcule $C_{th}^2$ et non $C_{th}$).}
\end{tcolorbox}

\vspace{4pt}

% -- QT2 -------------------------------------------------------------

\begin{tcolorbox}[questiontype,
  title={QT2 - Transmetteur 4-20\,mA - calcul du courant de sortie \hfill\normalfont\footnotesize E2 disciplinaire 2024 (partie E)}
]
\textbf{enonce.}
Un transmetteur de temperature a la plage : $T_{min} = -20^\circ\mathrm{C} \to 4\,\text{mA}$, $T_{max} = +200^\circ\mathrm{C} \to 20\,\text{mA}$.
Calculer $I_{out}$ pour $T = 120^\circ\mathrm{C}$.

\tcbline

\textbf{Corrige.}

\textbf{etape 1 - application de la relation lineaire.}
\[
  I_{out} = 4 + 16 \cdot \frac{T - T_{min}}{T_{max} - T_{min}}
          = 4 + 16 \cdot \frac{120 - (-20)}{200 - (-20)}
          = 4 + 16 \cdot \frac{140}{220}
\]

\textbf{etape 2 - resultat.}
\[
  I_{out} = 4 + 10{,}2 \approx \mathbf{14{,}2\,\text{mA}}
\]
\textit{Verif. :} $14{,}2\,\text{mA}$ est bien entre 4 et 20\,mA, coherent avec $T = 120^\circ\mathrm{C}$ dans la plage $[-20;200]^\circ\mathrm{C}$.\checkmark
\end{tcolorbox}

\vspace{4pt}

% -- QT3 -------------------------------------------------------------

\begin{tcolorbox}[questiontype,
  title={QT3 - Surveillance par PT100 - algorithme et interpolation \hfill\normalfont\footnotesize E2 disciplinaire 2022 (cogeneration)}
]
\textbf{enonce.}
Trois sondes PT100 ($T_1$, $T_2$, $T_3$) surveillent une canalisation vapeur - $P_s = 12\,\text{bars}$.
Lire la temperature de vaporisation dans la table et ecrire l'algorithme de surveillance.

\tcbline

\textbf{Corrige.}

\textbf{Lecture de table :} - $P_s = 12\,\text{bars}$ - $T_{vap} = 187{,}96^\circ\mathrm{C}$ (interpolation lineaire si necessaire).

\textbf{Algorithme de surveillance :}
\begin{center}
\renewcommand{\arraystretch}{1.2}
\begin{tabular}{l}
\textbf{SI} $(T_1 - T_2 < 12^\circ\mathrm{C})$ \textbf{ET} $(T_3 < 160^\circ\mathrm{C})$ \textbf{ALORS}\\
\quad etat $\leftarrow$ Fonctionnement correct\\
\textbf{SINON}\\
\quad etat $\leftarrow$ Alarme / Defaut\\
\textbf{FIN SI}
\end{tabular}
\end{center}
$(T_1 - T_2 < 12^\circ\mathrm{C})$ detecte un retour d'eau trop froid (condensats) ;
$T_3 < 160^\circ\mathrm{C}$ surveille la vapeur en fin de circuit.
\end{tcolorbox}

\end{tcolorbox}

% -----------------------------------------------------------------------
%  4. ATTENTES DU JURY
% -----------------------------------------------------------------------

\begin{tcolorbox}[
  jury,
  title={4. Attentes du jury - Erreurs recurrentes et points de vigilance},
]

\begin{enumerate}[label=\textbf{\arabic*.}, leftmargin=*, itemsep=5pt]

  \item \textbf{Oublier la racine carree dans $C_{th}$.}
  La formule donne $C_{th}^2$ : il faut imperativement extraire $\sqrt{\cdot}$ - la fin.
  Arreter au carre est l'erreur la plus signalee par le jury 2019.

  \item \textbf{Transmetteur 4-20\,mA : $\Delta I = 16\,\text{mA}$, pas 20.}
  Le courant varie de 4 - 20\,mA, soit une plage de $16\,\text{mA}$.
  Diviser par 20 au lieu de 16 est une erreur frequente.

  \item \textbf{Interpolation lineaire dans les tables thermodynamiques.}
  Si la pression exacte n'est pas tabulee, interpoler entre les deux valeurs encadrantes.
  Ne pas prendre la valeur la plus proche sans calcul.

  \item \textbf{Thermocouple : compensation de jonction froide.}
  La tension fournie est $e = k_S(T_{mesure} - T_{ref})$. Si $T_{ref} \neq 0^\circ\mathrm{C}$, appliquer la correction thermique avant de conclure.

  \item \textbf{PT100 vs thermocouple : ne pas confondre les principes.}
  PT100 : resistif (pont de Wheatstone) - grande precision, plage limitee.
  Thermocouple : effet Seebeck - plage etendue, sensibilite moindre.

  \item \textbf{Effet Peltier $\neq$ effet Seebeck.}
  Seebeck : $\Delta T \to e$ (capteur, thermocouple). Peltier : $I \to \Delta T$ (module thermoelectrique). L'effet exploite dans un thermocouple est l'effet Seebeck.

  \item \textbf{Preciser la plage et la tolerance du capteur choisi.}
  Le jury apprecie de verifier que la grandeur - mesurer est dans la plage du capteur et de preciser la classe de tolerance.

\end{enumerate}

\medskip
\begin{tcolorbox}[
  colback=colJuryFrame!12!colJury,
  colframe=colJuryFrame,
  left=4pt, right=4pt, top=2pt, bottom=2pt,
  fontupper=\small\itshape
]
Rapport jury 2019 : "Le calcul du couple thermique equivalent est bien connu en principe, mais les candidats font des erreurs dans la sommation : oublier certaines phases du cycle ou confondre $C_{th}^2$ et $C_{th}$."\\
Rapport jury 2024 : "Le jury apprecie les candidats qui precisent la plage de mesure et la tolerance du capteur choisi, et qui verifient la coherence entre la grandeur mesuree et la plage du capteur."
\end{tcolorbox}

\end{tcolorbox}

% -----------------------------------------------------------------------
%  PIED - PReSENCE DANS LES SUJETS
% -----------------------------------------------------------------------

\begin{tcolorbox}[
  enhanced, breakable,
  colback=colGrayLight,
  colframe=colSeparator,
  left=4pt, right=4pt, top=3pt, bottom=3pt,
  before skip=4pt, after skip=0pt,
  fonttitle=\small\bfseries,
  title={Presence dans les sujets analyses (2013-2026)},
]
\begin{center}
\footnotesize\renewcommand{\arraystretch}{1.15}
\begin{tabular}{lccccccccccccc}
\toprule
\textbf{Annee} & 13 & 14 & 15 & 16 & 17 & 18 & 19 & 20 & 21 & 22 & 23 & 24 & 25 \\
\midrule
E1 (commune)   & - & - & - & - & - & - & - & - & - & $\checkmark$ & - & $\checkmark$ & - \\
E2 discipl. IE & - & - & - & $\checkmark$ & - & $\checkmark$ & $\checkmark$ & $\checkmark$ & - & $\checkmark$ & - & $\checkmark$ & $\checkmark$ \\
\bottomrule
\end{tabular}
\end{center}
\smallskip
\footnotesize\textbf{Sous-themes frequents :}
couple thermique equivalent (5/13) - capteurs PT100/thermocouple (7/13) - transmetteur 4-20\,mA (4/13) - tables vapeur saturee (3/13) - resistance thermique composants (3/13)

\medskip
\textbf{Ordres de grandeur (reperes rapides).}
\begin{itemize}[leftmargin=1.3em]
  \item Reseau BT FR : tension phase-neutre typique \SI{230}{\volt}, tension phase-phase typique \SI{400}{\volt}, frequence \SI{50}{\hertz}.
  \item Courants d'etude frequents en sujet : de quelques \si{\ampere} a quelques dizaines de \si{\ampere}.
  \item En conversion d'energie, un rendement attendu reste generalement entre 0{,}8 et 0{,}98 selon la technologie.
\end{itemize}

\medskip
\textbf{Synthese express (a refaire sans support).}
\begin{enumerate}[leftmargin=1.4em]
  \item Poser les hypotheses du modele et les conventions de signe.
  \item Ecrire 2 a 4 relations cles avec unites.
  \item Realiser un calcul numerique complet et verifier l'ordre de grandeur.
  \item Justifier le resultat par un argument technique de niveau CAPET.
\end{enumerate}

\end{tcolorbox}

\CAPETENDFILE









